% А1.2 — Модуль: определение, геометрический смысл, свойства, базовые переходы
% Старая методичка: стр. 146–150 (старое задание 15)

\begin{tema}{А1.2}
\opr{Модуль} числа:\quad
$|a|=\begin{cases} a, & a\geqslant 0,\\ -a, & a<0.\end{cases}$\par
\opr{Смысл:} $|a|$ — расстояние от 0 до $a$, а $|a-b|$ — расстояние между точками $a$ и $b$ на прямой.

\begin{pfTablica}[bez shapki]{ll}
\pfPodshapka{Свойства}\\
$|a|\geqslant 0$, \ $\lvert -a\rvert=|a|$ & $|a|^2=a^2$, \ $\sqrt{a^2}=|a|$\\
$|ab|=|a|\,|b|$, \ $\left|\frac ab\right|=\frac{|a|}{|b|}$, $b\ne0$ & $|a-b|=|b-a|$\\
$|a+b|=|a|+|b| \iff ab\geqslant 0$ & $|a|+|b|=a+b \iff a,b\geqslant 0$\\
\pfPodshapka{Базовые переходы}\\
\Block[l]{1-2}{$|f|=c \iff f=\pm c$ \ при $c\geqslant0$; \ при $c<0$ решений нет}\\
$|f|<c \iff -c<f<c$ & $|f|>c \iff f>c$ \ или \ $f<-c$\\
\Block[l]{1-2}{\parbox{95mm}{\raggedright\footnotesize\color{pfSeryj}Для $\leqslant$ и $\geqslant$ — так же, с нестрогими знаками.
  Переходы для неравенств верны при любом $c$: при $c\leqslant0$ они сами дают «нет решений» или «все $x$».}}\\
\end{pfTablica}

\smallskip
\opr{Как раскрыть модуль:}
\begin{itemize}
  \item \emph{по определению} — два случая: $f\geqslant0$ и $f<0$;
  \item \emph{через расстояние} — когда модуль один и сравнивается с числом;
  \item \emph{по промежуткам} — когда модулей несколько: нули подмодульных
        выражений делят прямую, на каждом промежутке знаки известны.
\end{itemize}

\begin{primer}
$|x-2|\leqslant3$: точки на расстоянии не больше 3 от числа 2, то есть $x\in[-1;\,5]$.
\end{primer}
\risunok[0.5]{А1.2_1}
\end{tema}
